\scanpage{39}\mspage{38}
% L1–4
For a), it suffices to prove that $T_i$ is closed in \deleted{$\unsure{T|X}$} $A$ (the $T_j$, $j\ne i$, being so too, $T_i$ is also open \ldots\ ) If \deleted{\unsure{$T_i$ is f}} i.e. that
\[
\underbrace{\overline{T_i}}_{\inserted{\substack{\text{``closure'' in the}\\\text{tame sense}\\\text{of course!}}}}\,\cap\, A=T_{i},
\]
% L5–6
This is immediate

% L7–12
\uline{NB}\qquad If $T=T_i$ is a pseudo-segment piece, its tame ``closure'' is nothing other than the associated \inserted{closed} pseudo-segment piece ($=T$ if $T$ reduced to a poin\textsuperscript{t}) It suffices to see what happens when $T$ is not closed, at an \deleted{$a\in$} \deleted{$\unsure{\partial T}$} \deleted{$\mathrm{Fr}(T)=\dot T$ such that}
\[
a\in\partial T\setminus T\quad\ldots.
\]

% L13–14, b)
b) $T$ \uline{connected}\quad \inserted{If $T=\{\unsure{0|a}\}$ trivial. Otherwise $T$ a \deleted{\unsure{pseudo-}}segment piece.} \uline{NB}\quad $(\mathcal M_{L})_T=\mathcal M_T$,
% L15–17
regarding $T$ as a segment piece in it$\overline{\text{self}}$. \deleted{Then} Thus we may assume $T=L$. But If $X\subset\deleted{L}T$ tame,
% L18–19
then \unsure{:} \deleted{$\unsure{X}$} $X=X_{\Phi}=\bigcup_{i\in\unsure{I|F}}T_i$, then \deleted{Fr} the tame boundary is the union of the \deleted{$\unsure{\partial}$} $T_i^\bullet$,
% L20–22
it is empty if and only if the $T_i^\bullet$ are empty. $\checkmark$But $T_i^\bullet=\deleted{\text{empty}}\varnothing\Longrightarrow T_i=L$ so $X=\varnothing$ ($I=\varnothing$) or $X=L$ ($I\ne\varnothing$) OK.
